\documentclass[10.5pt,a4paper]{article}

\usepackage[utf8]{inputenc}
\usepackage[T1]{fontenc}
\usepackage{lmodern}
\usepackage{geometry}
\usepackage{amsmath, amssymb}
\usepackage{graphicx}
\usepackage{hyperref}

\geometry{margin=2.5cm}

\title{\textbf{Master Thesis Outline} \\ Loop Closure and Pose Graph Optimization for Sensorless Freehand 3D Ultrasound Reconstruction}
\author{Jannis Meyer}
\date{\today}

\begin{document}

\maketitle

\section{Background and Motivation}
Three-dimensional (3d) ultrasound (UR) allows for better and more precise diagnosis than simple 2d UR. As dedicated 3d UR scanners come with high costs, low convenience and lower image resolution compared to 2d UR, the concept of reconstructing 3d UR from freehand 2d UR images gained prominence. These are acquired by freehand-sweeping a standard 2d UR probe over the designated area. 3d UR reconstruction includes estimating the pose of the probe, acquiring images and with that doing the actual volumetric reconstruction. For pose estimation, the field has evolved from simply tracking the probe to using speckle correlation and finally to using deep learning methods (DL).\\
Current research focuses on solely feeding US images to DL models for pose estimation, without utilizing data from external sensors, tracking devices or Inertial Measurement Units (IMUs). The training data being used consists of 2d US images acquired from either only linear sweeps or linear, S-shaped and C-shaped sweeps in most cases. The most common error which is used for performance and comparison is Final Drift (FD), which is the reconstruction error of the final frame of a scan, created by the accumulated reconstruction errors over all frames. While state-of-the-art models showcase very low FD on the mentioned scan shapes, research on reconstruction quality for scans featuring frame redundancy, meaning circle shaped and back-and-forth shaped sweeps, is rare and includes using external sensors or IMUs together with DL for pose estimation.\\
This work aims to tackle this research gap by assessing 3d UR reconstruction quality using a combination of DL and global optimization techniques (Loop Closure and Pose Graph Estimation in particular), especially regarding the reconstruction of circular and back-and-forth sweeps. The main research question reads as:\\

\noindent\textit{How will global optimization techniques help improve the performance of DL-only based methods for pose estimation in the context of 3d UR reconstruction regarding scans featuring frame redundancy?}

\section{Setup}
As time is scarce, already existing data and pre-trained models will be used. This research might include fine-tuning or extra training of a model, if time allows it, but most of the work will happen at the inference stage.

\subsection{Model}
The model being used as base line and for further development is the one proposed in \href{https://github.com/ucl-candi/freehand}{this} repository. It uses an EfficientNet-B1 as a backbone with a fully connected layer at the end which feeds parameters to a volume reconstruction module. FD seems rather poor compared to other models (the papers mentioned in the repository are not clear on the quantitative performance of the public model, though images of example reconstructions show an approximate FD Rate, which is the FD divided by total scan length, of 15\%). This potential for improvement is important for this work as it gives room for comparison. 

\subsection{Data}
Two data sets were found fitting for this research, providing back-and-forth-sweeps and circular sweeps respectively:\\
\textbf{\href{https://zenodo.org/records/7740734}{Freehand\_US\_data}:} This data set includes twelve forearm scans from 19 volunteers, with a total scan length of 100-200mm. Straight, C-shaped and S-shaped sweeps were used to acquire the data. Redundancy in a scan may be artificially created by mirroring straight sweeps or duplicating parts of a scan. The data set was used for training and testing (split) the model mentioned before, which makes it suitable for comparison.\\
\textbf{\href{https://tdsc-abus2023.grand-challenge.org/Dataset/}{ABUS Data Set}}: This data set contains 200 volumes of breast tissue for breast cancer detection, acquired by a fixed, circular UR scanner. As the data does not stem from freehand UR scanning and serves the purpose of detection rather than reconstruction, measures have to be taken. Firstly, the data is 3d, so 2d data has to extracted by slicing up volumetric data and assigning pose information by adequate calculation. Furthermore, artificial noise has to be added to mimic freehand scanning, as the scans were performed by a fixed scanner. Data sets for 3d reconstruction specifically featuring circular scans could not be found, which is why the proposed data set has to suffice.\\
A similar alternatives to this is \href{https://ipeusctdb1.ipe.kit.edu/~usct/challenge/?page_id=183}{2D CSIC/UCM USCT}, for which a scanner moving around a tissue phantom was used. This data is 2d, meaning only artificial noise would have to be added.

\section{Methodology and Expected Outcome}
As stated, the main objective of this research is the performance enhancement from global optimization techniques for freehand 3d US reconstruction using DL only. For that, several methods shall be compared:
\begin{itemize}
\item Base Line: Raw proposed DL model on both data sets from last step
\item DL model + Pose Graph Optimization: Tested on both data sets
\item DL model + Pose Graph Optimization with Loop Closure: Tested on ABUS data set (circular motion only)
\item DL model + Optical Flow : Tested on both data sets (optional)
\end{itemize}

Statistics for comparing (e.g. averaged four corners of a frame as reference point):
\begin{itemize}
\item FD and FD Rate (especially at loop closure of circular data)
\item Average reconstruction error over all frames
\item Maximum reconstruction error of a scan over all frames
\item Volumetric error between reconstructed and ground truth volume
\end{itemize}

\section{Limitations and Concerns}
\begin{itemize}
\item Dependency on performance of open-source model and quality of publically available data
\item Artificial redundancy-featuring data
\item Global optimization techniques not increasing performance
\item Limited time
\end{itemize}

\bibliographystyle{plain}
\bibliography{references}

\end{document}